% 第 3 周周记（2026.08.17 - 08.21）
\documentclass[UTF8,fontset=fandol]{ctexart}
% =====================================================================
% internship/preamble.tex
% 实习周记统一模板的共用导言区。各周文件在 \documentclass 之后用
%   % =====================================================================
% internship/preamble.tex
% 实习周记统一模板的共用导言区。各周文件在 \documentclass 之后用
%   % =====================================================================
% internship/preamble.tex
% 实习周记统一模板的共用导言区。各周文件在 \documentclass 之后用
%   \input{../preamble.tex}
% 引入本文件，正文前用 \weekhead{周次}{时间}{主题} 输出抬头。
% 编译命令：xelatex weekN.tex
% =====================================================================

\usepackage[a4paper,top=2.6cm,bottom=2.6cm,left=2.8cm,right=2.8cm]{geometry}
\usepackage{booktabs}   % 表格横线，week3 的结果对比表会用到
\usepackage{array}

% 周记抬头：\weekhead{周次}{时间区间}{本周主题}
\newcommand{\weekhead}[3]{%
\begin{center}
{\zihao{3}\bfseries 实习周记（#1）}\\[10pt]
{\zihao{-4}姓名：【姓名】\qquad 学号：【学号】}\\[4pt]
{\zihao{-4}时间：#2}\\[4pt]
{\zihao{-4}主题：#3}
\end{center}
\vspace{2pt}
\hrule height 0.8pt
\vspace{10pt}
\zihao{-4}%
}

% 引入本文件，正文前用 \weekhead{周次}{时间}{主题} 输出抬头。
% 编译命令：xelatex weekN.tex
% =====================================================================

\usepackage[a4paper,top=2.6cm,bottom=2.6cm,left=2.8cm,right=2.8cm]{geometry}
\usepackage{booktabs}   % 表格横线，week3 的结果对比表会用到
\usepackage{array}

% 周记抬头：\weekhead{周次}{时间区间}{本周主题}
\newcommand{\weekhead}[3]{%
\begin{center}
{\zihao{3}\bfseries 实习周记（#1）}\\[10pt]
{\zihao{-4}姓名：【姓名】\qquad 学号：【学号】}\\[4pt]
{\zihao{-4}时间：#2}\\[4pt]
{\zihao{-4}主题：#3}
\end{center}
\vspace{2pt}
\hrule height 0.8pt
\vspace{10pt}
\zihao{-4}%
}

% 引入本文件，正文前用 \weekhead{周次}{时间}{主题} 输出抬头。
% 编译命令：xelatex weekN.tex
% =====================================================================

\usepackage[a4paper,top=2.6cm,bottom=2.6cm,left=2.8cm,right=2.8cm]{geometry}
\usepackage{booktabs}   % 表格横线，week3 的结果对比表会用到
\usepackage{array}

% 周记抬头：\weekhead{周次}{时间区间}{本周主题}
\newcommand{\weekhead}[3]{%
\begin{center}
{\zihao{3}\bfseries 实习周记（#1）}\\[10pt]
{\zihao{-4}姓名：【姓名】\qquad 学号：【学号】}\\[4pt]
{\zihao{-4}时间：#2}\\[4pt]
{\zihao{-4}主题：#3}
\end{center}
\vspace{2pt}
\hrule height 0.8pt
\vspace{10pt}
\zihao{-4}%
}


\begin{document}

\weekhead{第 3 周}{2026 年 8 月 17 日至 8 月 21 日}{离线增量管道实现与评估实验}

这一周基本都在写代码和跑实验。先搭了 trace 驱动的仿真器，把 Erigon 的块内语义复刻出来：每个块开始前清空缓存，槽位在块内的第一次访问计冷访问代价 3 微秒，之后的重复访问计热代价 30 纳秒，交易按块内顺序回放。仿真器先实现了两个对照组，E0 是不做任何预取的基线，E1 用交易真实的访问集合充当预测结果，作为理论上界。E1 的 recall 应该等于 1.0，对不上就说明仿真逻辑有 bug，我用它来自检。E0 在评估窗口的总代价是 1,128,660 微秒，E1 只有 14,649 微秒，相差约 77 倍，这个空间就是预取能争取的收益上限。

模型在上一阶段已经训练过，输入 8 个特征，包括目标合约地址、函数选择器、合约代码哈希、calldata 前三个参数、发送方地址和交易金额，标签是交易真实访问的槽位集合，用 LightGBM 多标签分类，候选槽取前 1000 个。本周把数据切分改成了按区块号时序 70/30 分成 training window 和 evaluation window，此前用的是随机切分，我自查时发现这种方式会把未来的信息漏进训练集，于是全部重跑。离线增量管道是方法的核心：在训练窗口内只对查表未命中的交易做预测，把预测结果和真实访问求交集，按 (pattern key, storage slot) 累计正确次数，再做最小支持度和每键新增上限两重过滤，最后合并回哈希表。候选实例累计 20 多万条，过滤后留下 253 个新键和 2,641 个槽位，平均每个键 10.4 个槽。

第一次跑出完整对比是在周三晚上。E2 使用原始哈希表，recall 为 0.2223，总存储访问代价 880,983 微秒；E3 使用离线增强后的表，recall 为 0.3117，总代价 781,405 微秒。recall 相对提升 40.2\%，代价下降 11.3\%。具体数字如下表。precision 从 0.1888 降到 0.1648，属于覆盖面扩大后的正常取舍，多出来的误预取代价小于多覆盖带来的收益。看到这组数字时觉得前几周的工夫有了交代，不过我也提醒自己，还要证明结果不是偶然。

\begin{table}[h]
\centering
\begin{tabular}{lcccc}
\toprule
实验组 & recall & precision & 总存储访问代价（微秒）\\
\midrule
E2：原始哈希表 & 0.2223 & 0.1888 & 880,983 \\
E3：离线增强 & 0.3117 & 0.1648 & 781,405 \\
\bottomrule
\end{tabular}
\end{table}

之后几天都在做验证。消融实验里，E0 到 E1 的差距代表理论改进空间，E2 填补了其中的 22.2\%，E3 进一步填补到 31.2\%。缩放实验把数据规模从 2K 笔、10K 笔一直跑到全量 94K 笔，recall 增益从 26.8\% 单调上升到 40.2\%，说明改进来自更多历史信号，而不是偶然。敏感性分析把冷热延迟参数各扰动正负 20\%，加速比波动只有约 0.1\%，recall 和 precision 只取决于预测是否正确，与延迟参数无关，这部分我用解析式验证。自查还翻出几组早期实验只在小样本上跑过，规模不够有说服力，这周全部重跑并把 CSV 按实验名存档。回头看，如果没有"数据留痕、口径固定"这套纪律，后面写论文时根本说不清每个数字的出处。

\end{document}
